\chapter{Literature Review}\thispagestyle{empty} %chapter head
\graphicspath{{Chapter2/figures/}}
{ \setstretch{1.5}
%--------------------------------------------------------------------------------------
The "Related Work" section of a project report is where you review and discuss existing research, projects, or methodologies that are relevant to your work. This section helps to contextualize your project within the broader field, demonstrating how it builds upon or differs from what has already been done.

\section{Steps to Write the Related Work Section:}
\subsection{Identify Relevant Work:}

Begin by identifying and selecting previous research papers, articles, projects, or methodologies that are relevant to your project. These can include academic papers, industry reports, books, or even well-known software or systems related to your topic.
\subsection{Organize the Literature:}

Group the related work into categories based on themes, methodologies, or approaches. This helps in creating a structured and coherent discussion. For example, you could categorize them by different approaches to solving a similar problem, by technology used, or by chronological development.
\subsection{Summarize and Compare:}

Summarize the key points of each piece of related work. Highlight the objectives, methods, findings, and limitations of each.
Compare and contrast these works with your project. Discuss how your project is similar to or different from these studies or systems. Explain how your work builds on, extends, or deviates from existing work.
\subsection{Highlight the Gaps:}

Identify gaps or limitations in the existing work that your project aims to address. This justifies the need for your project and positions it as a valuable contribution to the field.
\subsection{Cite Appropriately:}

Ensure you cite all the related work correctly using the appropriate citation style (APA, MLA, IEEE, etc.) as per the guidelines of your report.

\section{Example of a Related Work Section:}

The field of database security has been extensively studied, with various approaches proposed to safeguard sensitive data. Smith et al. (2015) developed an encryption-based approach to securing databases, which focuses on encrypting data at rest using AES encryption. While this method effectively protects data stored on disk, it has performance drawbacks, particularly in high-transaction environments, as identified by Johnson and Lee (2018), who proposed a hybrid encryption model that balances security and performance.

Another significant contribution is the role-based access control (RBAC) model discussed by Kumar and Patel (2016), which ensures that users only have access to the data necessary for their roles. This approach has been widely adopted in enterprise environments due to its scalability and ease of management. However, Martinez and Gupta (2019) pointed out that RBAC can become complex to manage in dynamic organizations where roles frequently change.

Recent work by Liu et al. (2020) introduced a machine learning-based anomaly detection system to identify potential security breaches in real-time. Their system shows promise in detecting unusual database activities but requires significant computational resources, which may not be feasible for smaller organizations.

Despite these advancements, there remains a gap in providing a comprehensive security framework that integrates encryption, access control, and real-time monitoring in a lightweight, efficient manner suitable for medium-sized healthcare institutions. This project addresses this gap by developing an integrated security framework tailored for XYZ Hospital's database system, combining encryption, role-based access control, and a streamlined anomaly detection mechanism.

\section{Tips for Writing the Related Work Section:}
\begin{enumerate}
	\item Be Concise: Avoid lengthy descriptions. Focus on summarizing the most relevant aspects of each work.
\item Stay Focused: Only include work that directly relates to your project. Irrelevant details can dilute the effectiveness of this section.
\item Be Critical: Don't just describe existing work; analyze it. Highlight strengths, weaknesses, and areas where your work differs.
\item Use Clear Transitions: Clearly indicate how each piece of related work connects to the next, and how they collectively relate to your project.

\end{enumerate}
By following these steps and tips, you can craft a well-organized and insightful "Related Work" section that effectively situates your project within the broader context of the field.
 
%--------------------------------------------------------------------------------------
}